\section*{Bab \ref{ch:b2:retrosynthesis} --- Retrosintesis dan Strategi Multitahap}

\begin{solution}{exo:b2:retrosynthesis:1}
\begin{itemize}
  \item 2-Fenilpropan-2-ol: propanon + \ce{PhMgBr}, atau asetofenon + \ce{CH3MgBr}.
  \item Pentan-3-ol: propanal + \ce{CH3CH2MgBr}.
  \item 1-Sikloheksiletanol: etanal + sikloheksilmagnesium bromida, atau
        sikloheksanakarbaldehida + \ce{CH3MgBr}.
  \item 2-Metilbutan-2-ol: propanon + \ce{CH3CH2MgBr}, atau butanon + \ce{CH3MgBr}.
  \item 1-Feniletanol: benzaldehida + \ce{CH3MgBr}, atau etanal + \ce{PhMgBr}.
\end{itemize}
\end{solution}

\begin{solution}{exo:b2:retrosynthesis:2}
\omterm{def:b2:retrosynthesis:disconnection}{Sinton} akseptor $\mathrm{Ph\overset{+}{C}H{-}OH}$, ekuivalennya benzaldehida;
\omterm{def:b2:retrosynthesis:disconnection}{sinton} donor \ce{CH3-}, ekuivalennya metilmagnesium bromida (atau
metillitium). (Potongan yang lain, dengan \ce{Ph-} sebagai donor, memakai
\ce{PhMgBr} dan etanal.)
\end{solution}

\begin{solution}{exo:b2:retrosynthesis:3}
$0.90 \times 0.85 \times 0.80 \times 0.95 \times 0.70 = 0.407$: sekitar 41~\%.
\end{solution}

\begin{solution}{exo:b2:retrosynthesis:4}
\ce{LiAlH4}, yang diperlukan untuk mereduksi ester menjadi alkohol, juga
akan mereduksi keton. Lindungi keton itu sebagai asetal siklik (etana-1,2-diol,
katalis asam, air dikeluarkan); reduksi esternya dengan
\ce{LiAlH4}; lepaskan asetalnya dengan asam berair encer: 5-hidroksipentan-2-on.
\end{solution}

\begin{solution}{exo:b2:retrosynthesis:5}
FGI: butana-1,3-diol $\Rightarrow$ 3-hidroksibutanal (reduksi aldehida,
\ce{NaBH4}). \omterm{def:b2:retrosynthesis:disconnection}{Pemutusan} dua gugus aldehida $\beta$-hidroksi ini pada ikatan
$\alpha$--$\beta$: dua molekul etanal (\omterm{def:b2:enolates-aldol:aldol}{aldol}). Arah maju: etanal, basa encer,
dingin; lalu \ce{NaBH4}.
\end{solution}

\begin{solution}{exo:b2:retrosynthesis:6}
Cincin sikloheksena dipotong pada kedua ikatan \ce{C-C} yang dibentuk dalam
\omterm{def:b2:frontier-orbitals:diels-alder}{reaksi Diels--Alder}, yaitu ikatan-ikatan alilik terhadap ikatan rangkap
cincin pada sisi substituen: buta-1,3-diena dan but-3-en-2-on, dengan gugus
asetil pada \omterm{def:b2:frontier-orbitals:diels-alder}{dienofil}.
\end{solution}

\begin{solution}{exo:b2:retrosynthesis:7}
Dengan penomoran C1 (\ce{C=O}), C2=C3, lalu C4 dengan kedua metil
(\cref{met:b2:conjugate-additions:robinson-disconnection}): potongan \omterm{def:b2:enolates-aldol:aldol}{aldol}
C2--C3 menyisakan aldehida pada C3 dan sebuah metil keton; potongan Michael
C4--C5 menyisakan 2-metilpropanal (C3 karbonilnya, C4 karbon
$\alpha$-nya) dan but-3-en-2-on.
\end{solution}

\begin{solution}{exo:b2:retrosynthesis:8}
Satu tahap: benzena dengan butanoil klorida dan \ce{AlCl3} (\omterm{def:b2:aromatic-substitution:friedel-crafts}{asilasi
Friedel--Crafts}, tanpa penataan ulang, asilasi tunggal). Dua tahap:
benzaldehida dengan propilmagnesium bromida, lalu oksidasi alkohol
sekundernya. Jalur Friedel--Crafts lebih pendek dan pereaksinya lebih murah,
tetapi memakai satu ekuivalen penuh aluminium klorida, yang berakhir dalam
limbah berair.
\end{solution}

\begin{solution}{exo:b2:retrosynthesis:9}
Piridinium klorokromat dalam diklorometana, oksidasi Swern, atau periodinana
Dess--Martin menghasilkan heksanal; dikromat atau permanganat berair yang
diasamkan menghasilkan asam heksanoat, karena dalam air aldehida membentuk
hidratnya, yang dioksidasi lagi (\cref{prop:b2:retrosynthesis:mild-oxidation}).
\end{solution}

\begin{solution}{exo:b2:retrosynthesis:10}
Amina bereaksi lebih dahulu, tanpa dilindungi, jika alkoholnya dilindungi:
lindungi alkohol sebagai eter silil (yang bertahan pada asilasi); asilasi
aminanya; lepaskan gugus silil dengan fluorida; asilasi alkoholnya. Jika
amina harus dibiarkan bebas sementara alkohol bereaksi, amina itu akan
dilindungi sebagai karbamat \textit{tert}-butoksikarbonil (dilepaskan oleh
asam), yang ortogonal terhadap eter silil (dilepaskan oleh fluorida):
masing-masing dapat dibebaskan tanpa menyentuh yang lain.
\end{solution}

\begin{solution}{exo:b2:retrosynthesis:11}
\ce{CH3COCH2CH2COCH3} adalah senyawa 1,4-dikarbonil: potong ikatan \ce{C-C}
tengahnya. Donornya adalah \omterm{def:b2:enolates-aldol:enolate}{enolat} etil 3-oksobutanoat (natrium etoksida);
akseptornya, sebuah karbon $\alpha$ keton, disediakan oleh
1-bromopropan-2-on, \ce{BrCH2COCH3} (substitusi bimolekular bromida). Lalu
hidrolisis berair ester dan pemanasan melepaskan \ce{CO2}: heksana-2,5-dion.
\end{solution}

\begin{solution}{exo:b2:retrosynthesis:12}
Retrosintesis: metil keton dengan \ce{CH2} di sebelahnya adalah produk
sintesis ester asetoasetat; potong ikatan antara C3 dan C4: \omterm{def:b2:enolates-aldol:enolate}{enolat} etil
3-oksobutanoat dan halida alilik 1-bromo-3-metilbut-2-ena,
\ce{(CH3)2C=CHCH2Br}. Arah maju: natrium etoksida dalam etanol, lalu
halidanya; lalu basa berair, pengasaman, dan pemanasan, yang menghidrolisis
ester dan melepaskan \ce{CO2}
(\cref{prop:b2:enolates-aldol:decarboxylation}): 6-metilhept-5-en-2-on.
\end{solution}

\begin{solution}{pb:b2:retrosynthesis:1}
\textbf{1.} \ce{HO-C6H4-CH2CH2COCH3}, dengan \ce{OH} para terhadap rantai: sebuah fenol dan sebuah keton.
\textbf{2.} Hidrogenasi sebuah \ce{C=C}: \ce{ArCH2CH2CO} $\Rightarrow$ \ce{ArCH=CHCO}, sebuah
keton $\alpha,\beta$-tak jenuh.
\textbf{3.} (\textit{E})-4-(4-hidroksifenil)but-3-en-2-on, \ce{HOC6H4CH=CHCOCH3}.
\textbf{4.} Karbonil $\alpha,\beta$-tak jenuh: \omterm{def:b2:enolates-aldol:aldol}{kondensasi aldol}, dipotong pada \ce{C=C}.
\textbf{5.} 4-Hidroksibenzaldehida dan propanon.
\textbf{6.} Aldehida itu tidak mempunyai hidrogen $\alpha$ dan hanya dapat
menjadi elektrofil; propanon, yang berlebih, adalah satu-satunya sumber
\omterm{def:b2:enolates-aldol:enolate}{enolat} dan bereaksi dengan aldehida yang lebih reaktif alih-alih dengan
dirinya sendiri.
\textbf{7.} (1) 4-hidroksibenzaldehida, propanon berlebih, natrium hidroksida
berair; pengasaman. (2)
\ce{H2} pada paladium di atas karbon, pada suhu kamar, sampai satu ekuivalen terserap.
\textbf{8.} Gugus \ce{CHO} menstabilkan fenoksida, sehingga fenol ini
sekurang-kurangnya sama asamnya dengan fenol
($\mathrm pK_a$ 9.99). Pada $\mathrm{pH} > 13$, $[\ce{ArO-}]/[\ce{ArOH}] > 10^{13-10} = 10^3$:
fenoksida.
\textbf{9.} \ce{O-} mendorong kerapatan elektron melalui cincin ke karbon
karbonil (donor resonansi yang para terhadapnya): aldehidanya kurang
elektrofilik, dan kondensasinya lebih lambat.
\textbf{10.} Eter benzil: hidrogenasi pada paladium yang mereduksi \ce{C=C}
juga memutus eter benzil (hidrogenolisis), sehingga membebaskan fenol pada
tahap yang sama.
\textbf{11.} Tiga: benzilasi, \omterm{def:b2:enolates-aldol:aldol}{kondensasi aldol}, hidrogenasi beserta
pelepasan perlindungan.
\textbf{12.} Tidak: dalam kondisi lunak (suhu kamar, paladium) cincin
\omterm{def:b2:huckel:aromatic}{aromatik} tidak tereduksi.
\textbf{13.} Tidak dilindungi: fenoksida hanya memperlambat kondensasi, yang
dapat diimbangi oleh waktu yang lebih lama atau pemanasan; satu tahap yang
dihemat lebih berharga daripada keuntungan laju.
\textbf{14.} 4-Iodofenol dan but-3-en-2-on.
\textbf{15.} \ce{HOC6H4I + CH2=CHCOCH3 + Et3N -> HOC6H4CH=CHCOCH3 + Et3NH+ + I-}, dengan katalis paladium.
\textbf{16.} Pada \ce{CH2} ujung, karbon $\beta$: gugus aril menyisip pada
ujung yang kurang terhalang, dan eliminasi $\beta$-hidrida lalu menghasilkan
enon (\textit{E}) yang terkonjugasi.
\textbf{17.} \omterm{def:b2:catalytic-cycles:oxidative-addition}{Adisi oksidatif} aril iodida pada \ce{Pd(0)}, koordinasi dan
insersi alkena, eliminasi $\beta$-hidrida, dan regenerasi \ce{Pd(0)} oleh
basa, yang mengambil \ce{HI}
(\cref{ch:b2:catalytic-cycles}).
\textbf{18.} $162.188/(220.005 + 70.091 + 101.193) = 162.188/391.289 \approx 41.4~\%$.
\textbf{19.} \ce{C7H6O2 + C3H6O -> C10H10O2 + H2O}: $162.188/(122.123 + 58.080) = 162.188/180.203
\approx 90.0~\%$.
\textbf{20.} 100~\%: setiap atom \ce{H2} dan enon berakhir dalam produk.
\textbf{21.} $164.204/(122.123 + 58.080 + 2.016) = 164.204/182.219 \approx 90.1~\%$.
\textbf{22.} \ce{C=C}: karena terkonjugasi dan tidak terhalang, ikatan itu
terikat dan terhidrogenasi pada paladium jauh lebih cepat daripada \ce{C=O}
keton; berhenti setelah satu ekuivalen hidrogen mempertahankan ketonnya.
\textbf{23.} Jalur \omterm{def:b2:enolates-aldol:aldol}{aldol}: aldehida, propanon, dan natrium hidroksida yang
murah, dengan air sebagai produk samping. Jalur Heck: aril iodida, katalis
paladium, dan but-3-en-2-on, yang sangat beracun.
\textbf{24.} Jalur \omterm{def:b2:enolates-aldol:aldol}{aldol}: dua tahap, pereaksi yang murah dan lebih aman,
ekonomi atom di atas 90~\%.
\textbf{25.} Tahap itu mengalikannya dengan 0.90: $0.782 \times 0.90 \approx 70~\%$.
\textbf{26.} $0.85 \times 0.92 = \boldsymbol{\approx 78~\%}$ secara keseluruhan.
\end{solution}
